% Terminal apparatus of the concise study: the scope note and the References
% for the sources this companion actually uses, at the loci it cites. The
% revision timestamp and the rights colophon follow from the shared
% generation record.

\runninghead{Scope and References}
\section{Appendix: Scope and Qualifications}

{\footnotesize
\textbf{Edition, Mass and occurrence.} The Mass treated is \latin{Dominica
vigesima quarta et ultima post Pentecosten}, II classis, of the \work{Missale
Romanum}, editio typica, Typis Polyglottis Vaticanis 1962, printed
pp.~422--424, nos.~1769--1778, as said on Sunday 22~November 2026, the
twenty-sixth Sunday after Pentecost by count, in the universal 1962 calendar.
Its orations and lessons are its own; its five chants are printed in the same
words as those of the Twenty-third Sunday after Pentecost (pp.~416--417). Its
place in 2026, the omission of the feast of St Cecilia, the colour, the
Gloria, the Credo and the Trinity Preface rest on the Missal's general
rubrics, on the note closing the Twenty-third Sunday (p.~417) and on the
Missal's printed direction after the Secret, as \texttt{research/context.md}
records. The Ordinary and the Office are not treated.

\textbf{Text and English.} The Latin of the ten texts was read on page images
of the CMAA facsimile of the 1962 typical edition and collated with the
Ratisbon printing of 1862 (Pustet); the collation is in
\texttt{propers/verified.md}, and the expansive study prints every text whole.
The English of the chants and orations is that of the Lasance \work{New Roman
Missal} (Benziger, 1945), public domain in the United States by non-renewal;
the lessons and the other Scripture are quoted in the Douay--Rheims in
Challoner's revision. None of this English is an approved liturgical
translation of the 1962 Missal. Where the Lasance English departs from the
Latin, single Latin words or short phrases are glossed; no English is composed
for a liturgical text. Chrysostom on Matthew and on Colossians is quoted in the
English of the Nicene and Post-Nicene Fathers, and Augustine on Psalm~129 in
the same series' abridged English, so that those quotations are of the
condensed English and not of his continuous text. Fromage is quoted in the
English of Dom Laurence Shepherd. The other Latin of the Fathers and medieval
writers is given with the English rendering made for the expansive study.

\textbf{Reception.} The three interpretations are developed in full in the
expansive study; their loci, editions and bounds are recorded in
\texttt{research/scope.md} and \texttt{research/interpretations.md}. Every
connection drawn between two texts of the Mass belongs to the interpretations
unless an author is named as making it; no Father cited here comments on this
Mass as a whole, and none of them speaks of the Eucharist at Mt 24:28. Gregory
is cited on Luke's form of the discourse, and Bede on Mark's form of it and on
Mark's verses of the Communion's saying. The Latin of Jerome was read on page
images; Augustine, Gregory and Bede in transcriptions of Migne not collated
with page images; Hilary, Aquinas on Matthew, Rupert and Sicard in
uncorrected optical text, whose quotations are normalised in spelling;
Aquinas on Colossians and on Psalm~43 in the Corpus Thomisticum text. Rupert
of Deutz, Sicard of Cremona and Lucien Fromage are ecclesiastical writers and
are cited in support only; Rupert and Sicard are cited for the chants alone.
No commentary on the Collect, Secret or Postcommunion survives among the
commentators read beyond Fromage's brief notes.

\textbf{Chronology.} Every biblical date on page~2 is generated from the
repository's Scripture chronology corpus; the expansive study's Scriptural
Date and Location appendix gives each passage's sources and evidence.

\textbf{Rights and use.} The Latin is published on 17 U.S.C. 103(b) together
with its public-domain witness, the Pustet printing of 1862. The Douay--Rheims
and the Lasance English are quoted from registered editions. The other
quotations are short and come from editions published before 1931, named
below, or from medieval and ancient Latin; the Corpus Thomisticum and
Wikisource texts are used only for Latin that is itself medieval or ancient.
Scripture, liturgical texts and quoted translations are not relicensed by the
repository's licence. This is a study aid to the 1962 Missal, not an official
liturgical book, translation or critical edition.
\par}

\section{References}

{\footnotesize
\begin{itemize}[leftmargin=1.2em,itemsep=.15em]
\item \work{Missale Romanum ex decreto Sacrosancti Concilii Tridentini
restitutum, Summorum Pontificum cura recognitum}, editio typica (Typis
Polyglottis Vaticanis, 1962), CMAA facsimile: \work{Rubricae generales};
Proprium de Tempore, pp.~416--417 and 422--424, nos.~1769--1778.
\item \work{Missale Romanum} (Ratisbon: Pustet, 1862), pp.~359--360.
\item Lasance, F.~X., and Walsh, F.~A., \work{The New Roman Missal in Latin and
English}, revised edition (New York: Benziger Brothers, 1945), pp.~736--740.
\item \work{The Holy Bible Translated from the Latin Vulgate} (Douay--Rheims,
Challoner revision), Project Gutenberg text: Gen 43:16, 34; 50:17--20; Ps
21:17; 23:6; 43:4--12, 23; 75:5; 84:2--5; 102:5; 129:1--7; Is 40:31; 49:25;
Jer 29:5--14; Hab 2:3; Mt 6:6; 7:7; 24:3--36; Mk 11:22--25; Lk 21:29--31;
Jn 12:24; Rom 7:23; Col 1:4--14; 1~Jn 2:15.
\item Augustine, \work{Enarrationes in Psalmos} 43.8--10 and 84.3--5,
Migne Latin text (Wikisource); 129.1 and 129.3, NPNF, first series, vol.~8
(abridged); \work{Quaestiones Evangeliorum} I.37--39 and I.42, Migne Latin
text (Wikisource).
\item Bede, \work{In Marci evangelium expositio} III, on 11:25, and IV, on
13:14--16 and 13:31, Migne Latin text (Wikisource).
\item Chrysostom, John, \work{Homilies on Matthew} 75.2, 76.1--3 and 77.1,
NPNF, first series, vol.~10; \work{Homilies on Colossians} 2, NPNF, first
series, vol.~13.
\item Gregory the Great, \work{Homiliae in Evangelia} I.1, nn.~1, 3--5, Migne
Latin text (Wikisource).
\item Hilary of Poitiers, \work{In Evangelium Matthaei commentarius} 25.3--8
and 26.1--3, PL~9 (Paris, 1844), cols.~1053--1058; \work{Tractatus super
Psalmos}, on Ps 129, PL~9, col.~718.
\item Jerome, \work{Commentarii in Matthaeum} IV, on 24:15--35, PL~26 (Paris,
1845), cols.~177--181; \work{In Hieremiam} V, on 29:10--14, PL~24 (Paris,
1865), col.~893.
\item Thomas Aquinas, \work{Super Evangelium S.~Matthaei lectura}, cap.~24, in
\work{Opera omnia}, vol.~3 (Venice: Bettinelli, 1745), pp.~306--310;
\work{Super epistolam ad Colossenses}, cap.~1, lect.~3, Corpus Thomisticum
text; \work{Super Psalmos}, on Ps 43, nn.~3--4, Corpus Thomisticum text.
\item Fromage, Lucien, continuing Prosper Guéranger, \work{The Liturgical
Year: The Time after Pentecost}, vol.~2 (English vol.~11), trans.\ Laurence
Shepherd (Dublin: Duffy, 1883), ``The Twenty-fourth and Last Sunday after
Pentecost,'' pp.~505, 509--510, 513.
\item Rupert of Deutz, \work{De divinis officiis} XII.23, PL~170,
cols.~329--330.
\item Sicard of Cremona, \work{Mitrale} VIII.23 and VIII.25, PL~213.
\end{itemize}
\par}
